\begin{textsample}{NEG Dim 2 – 2013\_c6 – Score -2.00 – 2013\_c6/enade\_2013\_med\_vet\_5d625e2f60.txt}  \label{ex:f2_neg_057}
CAPÍTULO 3 ANÁLISE TÉCNICA DA PROVA Este capítulo tem por objetivo apresentar o desempenho dos estudantes concluintes de Medicina Veterinária no ENADE/2013.
Para isso, foram calculadas as estatísticas básicas da prova em seu todo ( seção 3.1.1 ), bem como as estatísticas dos componentes relacionadas ao Componente de Formação Geral ( seção 3.1.2 ) e ao de Conhecimento Específico da Área ( seção 3.1.3 ).
Dadas as suas características, foram analisadas em separado, as questões objetivas ( seção 3.2 ) e as questões discursivas ( seção 3.3 ).
Para as questões objetivas, foram disponibilizados os índices de facilidade e o de discriminação bisserial, também em separado para os Componentes de Formação Geral ( seção 3.2.1 ) e de Conhecimento Específico ( seção 3.2.2 ).
De cada componente, uma das questões foi escolhida para exemplificar a análise gráfica relacionando as alternativas escolhidas pelos estudantes ( inclusive o gabarito ) com o número de acertos no componente.
O Anexo I apresenta a íntegra da análise gráfica para todas as questões objetivas.
Para cada uma das questões discursivas, os conteúdos dos tipos mais comuns de respostas dos alunos são apresentados e comparados com o padrão de resposta esperado ( ver Anexo VIII com o padrão de respostas ).
Tomando como base as duas questões discursivas do Componente de Formação Geral, a seção 3.3.1.6 apresenta comentários sobre a correção das respostas com respeito à Língua Portuguesa.
Nas tabelas, são apresentados o tamanho da população inscrita e de presentes, e as seguintes estatísticas das notas11 : média do desempenho na prova, erro padrão da média, desvio padrão, nota mínima, mediana e nota máxima.
Tais estatísticas contemplam o total de estudantes concluintes da Área de Medicina Veterinária inscritos e presentes à prova do ENADE/2013, tendo -se em vista as seguintes agregações : ( a ) as Grandes Regiões e o país como um todo ; ( b ) a Categoria Administrativa ; e ( c ) a Organização Acadêmica.
Em relação aos gráficos de distribuição de notas, o intervalo considerado foi de 10 unidades, aberto à esquerda e fechado à direita, com exceção do primeiro intervalo, [ 0 ; 10 ], fechado em ambos os extremos.
Para os gráficos de distribuição das notas das questões discursivas, foram consideradas mais duas categorias : questão em branco e nota zero. 11 Essas estatísticas e outras estão definidas no Capítulo 1. 3.1 ESTATÍSTICAS BÁSICAS DA PROVA 3.1.1 Estatísticas Básicas Gerais A Tabela 3.1 apresenta as estatísticas básicas da prova por grande Região.
A população total de inscritos foi de 7.841.
Destes, 7.269 estiveram presentes, sendo 7,3 \% o índice de não comparecimento.
A Região de maior abstenção foi a Sudeste ( 9,2 \% ) e a de menor abstenção foi a Sul ( 4,0 \% ).
A média das notas da prova como um todo ( nas seções seguintes serão analisados os Componentes de Formação Geral e de Conhecimento Específico ) foi 45,3, sendo que os alunos da região Norte obtiveram a média mais baixa ( 42,2 ) e os da região Sul obtiveram a média mais alta ( 46,5 ).
As demais médias foram : 44,9 na região Nordeste, 45,5 na região Sudeste e 43,6 na região Centro-Oeste.
O desvio padrão para o Brasil como um todo foi 11,8, sendo o maior desvio padrão encontrado na região Centro-Oeste ( 12,7 ) e o menor nas regiões Norte e Sul ( 11,0 ), indicando uma menor dispersão das notas destas últimas regiões.
A região que obteve a maior nota máxima foi a Sudeste ( 81,1 ), ao passo que a região que atingiu a menor nota máxima foi a Norte ( 70,8 ).
A mediana do Brasil como um todo foi 45,6, sendo a maior mediana obtida na região Sul ( 46,8 ) e a menor obtida na Norte ( 42,0 ).
A nota mínima foi zero em todas as regiões, exceto na região Norte ( 10,2 ).
O comportamento das notas dos estudantes de todo o Brasil pode ser observado no Gráfico 3.1 que apresenta um histograma com a distribuição das mesmas.
Essa é uma distribuição unimodal com moda no intervalo ( 40 ; 50 ].
O coeficiente de assimetria da distribuição das notas é negativo ( -0,26 ) e não podemos considerar que a distribuição é aproximadamente simétrica.
As distribuições por Grande Região também apresentam assimetria negativa, concentração pouco maior do lado esquerdo do histograma e mais espalhada do lado direito, excetuando -se apenas a região Norte que possui assimetria positiva. espalhada do lado direito, excetuando -se apenas a região Norte que possui assimetria positiva.
Os Gráficos 3.2, 3.3 e 3.4 apresentam informações referentes à média da nota final dos participantes, desagregando os resultados de acordo com, respectivamente, as Grandes Regiões do país, a Categoria Administrativa e a Organização Acadêmica.
Os gráficos apresentam o valor da média das notas como uma barra e os extremos do intervalo de confiança de 95 \% como linhas verticais unidas por uma linha horizontal na forma da letra H maiúscula.
Considerando -se o gráfico de notas segundo Grande Região ( Gráfico 3.2 ), observa -se que existe diferença estatisticamente significativa ao nível de 95 \% entre a maior média, obtida na região Sul ( 46,5 ) e a menor, obtida na região Norte ( 42,2 ).
Levando -se em conta os estudantes por Categorias Administrativas ( Gráfico 3.3 ), observa -se que existe diferença estatisticamente significante entre as médias das notas das IES Públicas e Privadas.
Pode ser observado ainda que a média dos alunos de IES Públicas ( 49,1 ) é maior que a dos alunos de IES Privadas ( 42,9 ).
A diferença entre as médias das regiões Sul e Norte ( 4,3 ), a maior e a menor média, é inferior a diferença entre IES Públicas e Privadas ( 6,2 ).
Tendo como foco o Gráfico 3.4, que apresenta as notas médias das provas segundo Organização Acadêmica, constata -se que a diferença da média das notas dos estudantes provenientes de Universidades é estatisticamente significativa ao nível de 95 \% das médias das notas dos estudantes provenientes de Centros Universitários e Faculdades.
Já a diferença entre a média dos estudantes de Centros Universitários e Faculdades não é significativa estatisticamente.
A maior média foi obtida pelos estudantes de Universidades ( 46,9 ), e a menor, pelos de Faculdades ( 41,9 ). 3.1.2 Estatísticas Básicas no Componente de Formação Geral A Tabela 3.2 apresenta as estatísticas básicas em relação ao componente da prova que avalia a Formação Geral dos estudantes concluintes.
Os alunos de todo Brasil obtiveram desempenho médio de 48,0.
Quanto à variabilidade, o desvio padrão das notas dos estudantes do Brasil como um todo foi 15,1.
A maior média foi obtida na região Sul ( 49,0 ), e a menor, na região Norte ( 44,6 ).
As demais médias foram : 47,7 na região Nordeste, 48,3 na região Sudeste e 45,9 na região Centro-Oeste.
Já o maior desvio padrão foi obtido na região Centro-Oeste ( 16,3 ) e o menor na região Sul ( 14,0 ).
Os demais desvios padrões foram : 15,2 na região Norte, 16,2 na região Nordeste e 15,1 na região Sudeste.
A maior nota no Componente de Formação Geral da prova do ENADE foi obtida por pelo menos um aluno da região Norte ( 92,4 ) enquanto que a menor nota máxima foi obtida na região Centro-Oeste ( 81,5 ).
Nas outras regiões as notas máximas foram : 87,2 na região Nordeste, 87,6 na região Sudeste e 88,4 na região Sul.
A mediana do Brasil como um todo foi 49,1, sendo a menor mediana encontrada na região Norte ( 45,4 ) e a maior encontrada na A nota mínima nesta parte foi zero em todas as regiões, exceto na região Norte ( 7,5 ).
O Gráfico 3.5 propicia a avaliação do desempenho dos estudantes no Componente de Formação Geral a partir do histograma da distribuição das notas correspondentes.
A distribuição é unimodal, com moda em ( 50 ; 60 ], enquanto na prova como um todo a moda foi alcançada no intervalo ( 40 ; 50 ].
Nota -se, ainda, que no Gráfico 3.5 as notas apresentam uma maior dispersão do que no Gráfico 3.1 ( distribuição das notas da prova ), confirmado pela comparação dos desvios padrões : 11,8 para a nota da prova como um todo e 15,1 para o Componente de Formação Geral.
Para o Componente de Formação Geral, o coeficiente de assimetria da distribuição das notas dos estudantes é negativo ( -0,36 ).
A distribuição apresenta uma leve concentração à direita e cauda maior à esquerda.
Em todas as Grandes Regiões os histogramas também possuem assimetria negativa ( entre -0,55 para o Nordeste e -0,22 para o Norte ).
Nos Gráficos 3.6, 3.7 e 3.8 são apresentadas as informações referentes ao desempenho dos concluintes no Componente de Formação Geral, em diferentes agregações : Grande Região do país, Categoria Administrativa e Organização Acadêmica.
Observa -se pelo Gráfico 3.6 que existe diferença estatisticamente significativa ao nível de 95 \% entre a maior média das notas na Componente de Formação Geral, obtida na região Sul ( 49,0 ) e a menor, obtida na região norte ( 44,6 ).
Vemos que o intervalo de confiança mais largo é o da região Norte ; já o intervalo mais estreito é observado na região Sudeste.
Este fato está relacionado, também, com o tamanho da população envolvida, menor na região Norte do que na Sudeste.
No Gráfico 3.7, que representa as notas médias no Componente de Formação Geral segundo Categoria Administrativa do país, observa -se que existe diferença estatisticamente significativa entre as médias.
Assim como ocorreu para as notas da prova como um todo, para o Componente de Formação Geral, os concluintes de Medicina Veterinária, das IES Públicas ( 50,9 ) obtiveram uma média maior do que os das IES Privadas ( 46,1 ).
Considerando -se o tipo de Organização Acadêmica, nota -se, no Gráfico 3.8, que existe uma diferença estatisticamente significativa entre a maior e a menor média.
Nas Universidades, a média ( 49,4 ) é mais elevada do que em Centros Universitários ( 45,3 ) e Faculdades ( 44,9 ). 3.1.3 Estatísticas Básicas do Componente de Conhecimento Específico A Tabela 3.3 apresenta as estatísticas básicas referentes ao Componente de Conhecimento Específico da área de Medicina Veterinária.
A média do desempenho dos alunos do Brasil como um todo foi 44,4.
A maior média foi obtida na região Sul ( 45,7 ), e a menor, na região Norte ( 41,3 ).
As demais médias foram : 44,0 na região Nordeste, 44,6 na região Sudeste e 42,8 na região Centro-Oeste.
Quanto à variabilidade das notas, o desvio padrão do Brasil como um todo foi 12,9, sendo o maior desvio padrão observado na região Centro-Oeste ( 13,7 ) e o menor nas regiões Norte e Sul ( 12,3 ).
O desvio padrão da região Nordeste foi 13,3 e o da região Sudeste foi 12,9.
A mediana das notas dos estudantes de todo o Brasil foi 44,7.
A maior mediana ocorreu na região Sul ( 45,9 ) e a menor na região Norte ( 41,1 ).
As demais medianas foram 44,8 na região Nordeste, 44,6 na região Sudeste e 43,6 na região Centro-Oeste.
A nota máxima do Brasil como um todo foi 87,3, sendo obtida por pelo menos um aluno da região Nordeste.
As demais notas máximas foram : 76,2 na região Norte, 85,4 na região Sudeste, 83,7 na região Sul e 81,2 na região Centro-Oeste.
A nota mínima foi zero em todas as regiões, exceto na região Norte ( 8,1 ).
Assim como os Gráficos 3.1 e 3.5, o Gráfico 3.9, apresentado a seguir, proporciona uma avaliação do desempenho de concluintes em relação ao Componente de Conhecimento Específico com um histograma da distribuição das notas correspondentes.
Esta também é uma distribuição unimodal, e o grupo modal é o ( 40 ; 50 ], abaixo do grupo modal para a Formação Geral.
O coeficiente de assimetria da distribuição das notas do Componente de Conhecimento Específico é negativo ( -0,17 ).
Nota -se pelo histograma ( Gráfico 3.9 ) que esta distribuição apresenta uma leve concentração à direita e cauda maior à esquerda.
Em quase todas as Grandes Regiões os histogramas também possuem assimetria negativa ( entre -0,53 para o Nordeste e -0,07 para o Sudeste ), exceto o da região Norte ( 0,12 ) que possui assimetria positiva.
Os Gráficos 3.10, 3.11 e 3.12 apresentam, respectivamente, uma comparação dos resultados em relação à Grande Região do país, à Categoria Administrativa e à Organização Acadêmica, agora levando em conta o desempenho de alunos no Componente de Conhecimento Específico da prova.
Pelo Gráfico 3.10, observa -se que existe diferença estatisticamente significativa entre a maior e a menor média das notas no Componente de Conhecimento Específico, obtidas nas regiões Sul ( 45,7 ) e Norte ( 41,3 ) respectivamente.
Quanto à Categoria Administrativa ( Gráfico 3.11 ), observa -se um comportamento semelhante àquele da Parte de Formação Geral e à prova como um todo, ou seja, existe diferença estatisticamente significativa entre as médias das IES Públicas e Privadas, sendo que a maior média foi obtida por alunos de IES Públicas de ensino ( 48,5 ).
Quanto ao Gráfico 3.12, observa -se, mais uma vez, que existe diferença estatisticamente significativa ao nível de 95 \% entre as notas no Componente de Conhecimento Específico de Universidades e dos demais tipos de Organização Acadêmica.
Sendo que a média dos concluintes das Universidades ( 46,1 ) foi maior do que de Centros Universitários ( 41,3 ) e de Faculdades ( 40,9 ). 3.2 ANÁLISE DAS QUESTÕES OBJETIVAS 3.2.1 Componente de Formação Geral A Tabela 3.4 apresenta as estatísticas básicas relativas às oito questões objetivas do componente da prova que abrange a Formação Geral dos estudantes.
A média do Brasil foi 51,3.
A menor média foi encontrada na região Centro-Oeste ( 49,9 ) e a maior na região Nordeste ( 52,5 ).
As demais médias foram 50,7 na região Norte, 51,1 na região Sudeste e 51,8 na região Sul.
O desvio padrão do Brasil foi 18,3, sendo o maior desvio padrão encontrado na região Centro-Oeste ( 18,7 ) e o menor nas regiões Norte e Sul ( 17,8 ).
Os demais desvios foram : 18,5 na região Nordeste e 18,3 na região Sudeste.
As mesmas medianas ( 50,0 ), notas máximas ( 100,0 ) e notas mínimas ( 0,0 ), nas questões objetivas do Componente de Formação Geral, foram obtidas em todas as regiões, sem exceção.
A Tabela 3.5 apresenta o índice de facilidade e o índice de discriminação ( ponto bisserial ) para cada uma das questões objetivas do Componente de Formação Geral.
Quanto ao índice de facilidade, foram usadas as seguintes cores para diferenciar o nível de dificuldade da questão : • Azul para as questões classificadas com índice muito fácil ( índice \&gt ; =0,86 ), verde para as questões classificadas com índice fácil ( 0,61 a 0,85 ), amarelo para as questões classificadas com médio ( 0,41 a 0,60 ), vermelho para as questões classificadas com difícil ( 0,16 a 0,40 ) e roxo para as questões classificadas com muito difícil ( \&lt ; =0,15 ).
Já quanto ao índice de discriminação, foram usadas as seguintes cores para qualificar a questão : • As questões classificadas com índice fraco receberam a cor vermelha ( índice \&lt ; =0,19 ), as classificadas com médio receberam a cor amarelo ( 0,20 a 0,29 ), as classificadas com bom receberam a cor verde ( 0,30 a 0,39 ) e as classificadas com muito bom ( \&gt ; =0,40 ) receberam a cor azul.
As questões objetivas do Componente de Formação Geral, segundo o índice de facilidade, foram assim avaliadas : das oito questões, uma teve o índice de facilidade classificado como muito fácil e duas como fácil.
Três questões foram tidas como médio, por terem índice de acertos situado na faixa entre 0,41 e 0,60 ( de 41,0 \% a 60,0 \% de acertos ).
Duas questões foram consideradas de dificuldade difícil, situando -se no intervalo entre 0,16 e 0,40 do índice de facilidade, ou seja, houve entre 16,0 \% e 40,0 \% de acertos.
Como já comentado, para análise das questões objetivas relativas à Formação Geral segundo o poder de discriminação, utilizou -se, o índice de discriminação ponto bisserial.
Nesta análise, as questões foram assim avaliadas : cinco das oito questões apresentaram O Gráfico 3.13, para exemplificar, analisa o comportamento da questão de número 3 de Formação Geral.
Trata -se da quinta questão mais fácil e a que obteve o segundo maior índice de discriminação dessa parte da prova.
Neste gráfico, cada uma das cinco curvas representa o percentual de respostas em determinada alternativa da questão, em função da nota dos estudantes nesta parte da prova ( Formação Geral/Múltipla Escolha ), antes de possíveis eliminações pelo critério do ponto bisserial.
A curva em vermelho corresponde à alternativa E, a correta para esta questão.
Assim, observa -se que entre os estudantes com notas mais baixas, nessa parte do exame, a situação mais frequente foi a escolha de uma das alternativas incorretas : a alternativa A índices acima de 0,40 e, assim, foram classificadas com índice muito bom para esse grupo de alunos ; três questões tiveram índice de discriminação bom, entre 0,30 e 0,39.
Nenhuma questão foi classificada com um índice fraco e, portanto, nenhuma questão foi eliminada pelo critério ponto bisserial.
O índice de facilidade variou de 0,24 a 0,96, e o de discriminação, de 0,30 a 0,49.
As questões com índice de discriminação muito bom, de números 3, 4, 6, 7 e 8, figuraram entre as mais fáceis desse conjunto : duas ( questões 4 e 8 ) classificadas na categoria fácil e outras três na categoria médio do índice de facilidade.
Em particular, a questão 8 apresentou maior poder discriminatório, com índice 0,49, e foi a segunda mais fácil, com uma proporção de 0,69 de acertos.
As questões de número 2 e 5 apresentaram índice de facilidade 0,24, ou seja, um quantitativo de 24,0 \% de estudantes conseguiu resolvê -las, dentro do universo de participantes.
Além disso, seu índice de discriminação foi bom em ambas, com valores iguais a 0,31 e 0,37 respectivamente.
O Gráfico 3.13, para exemplificar, analisa o comportamento da questão de número 3 de Formação Geral.
Trata -se da quinta questão mais fácil e a que obteve o segundo maior índice de discriminação dessa parte da prova.
Neste gráfico, cada uma das cinco curvas representa o percentual de respostas em determinada alternativa da questão, em função da nota dos estudantes nesta parte da prova ( Formação Geral/Múltipla Escolha ), antes de possíveis eliminações pelo critério do ponto bisserial.
A curva em vermelho corresponde à alternativa E, a correta para esta questão.
Assim, observa -se que entre os estudantes com notas mais baixas, nessa parte do exame, a situação mais frequente foi a escolha de uma das alternativas incorretas : a alternativa A ( em azul ) ou C ( em preto ).
Na medida em que a nota aumenta, indicando desempenho melhor nesta parte da prova, aumenta concomitantemente a proporção de estudantes que selecionaram a alternativa correta E, atingindo 100 \% para os alunos com 8 acertos.
Essa análise permite verificar como a questão discriminou os grupos de desempenho, justificando o alto índice obtido na questão.
Os gráficos relativos às demais questões de Formação Geral constam do Anexo I. 3.2.2 Componente de Conhecimento Específico A Tabela 3.6 apresenta as estatísticas básicas em relação às questões objetivas do Componente de Conhecimento Específico da prova, por Grande Região.
A média do Brasil deste componente foi de 48,9.
A menor média foi observada na região Norte ( 46,4 ) e a maior na região Sul ( 50,1 ).
O desvio padrão de todo o Brasil foi 14,1, sendo o menor desvio padrão encontrado na região Sul ( 13,4 ) e o maior na região Centro-Oeste ( 15,0 ).
A mediana de todo o Brasil foi 47,8.
Nas regiões Norte, Sudeste e Centro-Oeste, a mediana foi também 47,8.
Na região Nordeste e na Sul, a mediana foi 52,2.
A maior nota máxima da prova foi obtida, nas questões objetivas do Componente de Conhecimento Específico, por pelo menos um aluno na região Nordeste ( 95,7 ), enquanto a menor nota mínima foi obtida na região Norte ( 82,6 ).
As demais notas máximas foram : 91,3 na região Sudeste, 87,0 nas regiões Sul e Centro-Oeste.
A nota mínima foi zero em todas as regiões, exceto na região Norte ( 8,7 ).
A Tabela 3.7 apresenta o índice de facilidade e o índice de discriminação ( ponto bisserial ) para cada uma das questões objetivas do Componente de Conhecimento Específico.
Para facilitar a diferenciação das questões usaram -se as mesmas cores da Tabela 3.5 para as diferentes classificações dos índices de facilidade e de discriminação.
Dentre as questões objetivas da parte da prova relativa ao Componente de Conhecimento Específico, nenhuma foi anulada pela Comissão.
Desse modo, a classificação quanto ao índice de facilidade foi estabelecida com base nas 27 questões.
A partir dos índices obtidos, pode -se concluir que grande parte das 27 questões objetivas da prova foi considerada como difícil, 11 ao todo.
As 16 questões restantes foram classificadas quanto ao índice de facilidade da seguinte maneira : uma como muito fácil, quatro como fácil, dez como médio e uma como muito difícil.
Já quanto aos índices de discriminação das questões objetivas do Componente de Conhecimento Específico da prova, tem -se como resultado a seguinte classificação : 12 das 27 questões foram consideradas como bom, enquanto nenhuma teve índice de discriminação muito bom.
Assim, para 12 das 27 questões, os índices de discriminação foram bom ou muito bom.
Dentre as demais, 11 delas foram classificadas como médio e outras quatro como fraco, sendo 15, por conseguinte, a quantidade de questões nos dois patamares mais baixos de discriminação.
Constata -se, assim, que a prova – no que se refere ao Componente de Conhecimento Específico – possuía capacidade mediana de discriminar entre aqueles que dominam ou não o conteúdo.
Dentre as questões que alcançaram os maiores índices de discriminação, as de números 14, 17, 18, 21, 23, 24, 25, 28, 29, 30, 32 e 34 classificadas com índice bom, situando -se no intervalo de 0,30 a 0,39 do índice, uma delas ( questão 14 ) foi classificada na categoria muito fácil, quatro ( questões 17, 21, 24 e 34 ) como fácil, quatro ( questões 23, 28, 29 e 32 ) como médio e outras três como difícil, quanto ao índice de facilidade.
A questão de número 19 foi a mais difícil dentre as 27 questões específicas, com baixo índice de facilidade, apenas 9,0 \% de acertos.
Essa questão apresentou poder discriminatório igualmente baixo, 0,10, o que comprova ter sido esta a mais difícil para os estudantes.
Destaca -se, também, a questão 27, com índice de facilidade 0,18, o que, em termos percentuais, corresponde a 18,0 \% de estudantes que responderam acertadamente, obtendo, ainda, 0,18 de índice de discriminação.
Tais questões foram, portanto, duas das mais difíceis da prova.
As questões 19 e 27 foram eliminadas do cômputo da nota final pelo critério do ponto bisserial.
Além destas duas, as demais questões com índice fraco de discriminação, questões 15 e 20 também não foram computadas.
A título de exemplo das análises do comportamento das questões objetivas, o Gráfico 3.14 analisa a questão 14 do Componente de Conhecimento Específico.
Esta foi uma das questões classificadas como muito fácil nesta prova, apresentando índice de facilidade 0,87, ou seja, 87,0 \% dos estudantes assinalaram acertadamente a opção D, correspondente ao gabarito.
Seu índice de discriminação foi igual a 0,31, classificado como bom.
Neste gráfico, cada uma das cinco curvas representa o percentual de respostas em determinada alternativa da questão 14, em função da nota dos estudantes nesta parte da prova, antes de possíveis eliminações de questões pelo critério do ponto bisserial.
A alternativa correta D, representada no gráfico pela curva em roxo, foi escolhida em maiores proporções pelos alunos com desempenho melhor nesta parte da prova.
Já as alternativas incorretas, também denominadas distratores, foram selecionadas principalmente por aqueles com notas mais baixas.
A proporção de alunos que selecionou a resposta correta D aumenta gradativamente, chegando a atingir 100 \% para aqueles que acertaram pelo menos 22 questões, enquanto a proporção dos que escolheram alternativas incorretas decai, a partir da primeira nota não nula, como função da nota nesta parte da prova.
Os gráficos relativos às demais questões do Componente de Conhecimento Específico constam do Anexo I. 3.3 ANÁLISE DAS QUESTÕES DISCURSIVAS 3.3.1 Componente de Formação Geral As análises dos resultados de desempenho dos estudantes de Medicina Veterinária nas duas questões discursivas relativas à Formação Geral encontram -se na Tabela 3.8 e no Gráfico 3.15.
Na Tabela 3.8 observa -se que as notas médias foram mais baixas nesse conjunto de questões do que no das objetivas.
Os estudantes de todo o Brasil obtiveram, em Formação Geral, média 51,3 nas questões objetivas e 42,9 nas questões discursivas.
A mediana também confirma o pior desempenho dos alunos de todo o Brasil nas questões discursivas do Componente de Formação Geral.
Enquanto essa estatística foi de 48,0 para questões discursivas, para as questões objetivas a mediana foi de 50,0.
Pode -se, também, notar um aumento do desvio padrão de 18,3, nas questões objetivas do Componente de Formação Geral dos alunos de todo o Brasil, para 20,9 nas questões discursivas do mesmo componente.
Como já comentado, a mediana de todo o Brasil, neste componente, foi 48,0, sendo a maior mediana encontrada na região Sudeste ( 49,5 ) e a menor na região Norte ( 37,5 ).
A maior nota máxima foi obtida, nas questões discursivas do Componente de Formação Geral, por pelo menos um aluno da região Sul ( 90,5 ), enquanto a menor nota máxima foi obtida na região Norte ( 82,0 ).
Em todas as regiões a nota mínima foi zero, sem exceção.
O Gráfico 3.15 representa a distribuição das notas nas questões discursivas no Componente de Formação Geral.
A moda desta distribuição ocorre no intervalo ( 50 ; 60 ].
Devido à grande quantidade de notas zero e à alta frequência de alunos que deixaram este tipo de questão em branco o intervalo [ 0 ; 10 ] se apresenta como uma segunda moda.
A distribuição possui assimetria à esquerda, coeficiente de assimetria -0,76.
Em todas as regiões o coeficiente de assimetria também é negativo, o maior em módulo, -0,86, na região Sul, e o menor na região Norte, -0,17.
A avaliação das questões discursivas de Formação Geral considerou em separado o conteúdo ( peso 0,8 na nota ) e o desempenho em língua portuguesa ( peso 0,2 ).
Na sequência, os resultados verificados para cada uma das questões discursivas de Formação Geral com relação ao conteúdo serão apresentados, estabelecendo -se relações com os temas abordados em cada uma delas.
Os comentários da Banca de docentes corretores a respeito do observado na correção das respostas dos estudantes, suas impressões e conclusões serão apresentados junto à análise de cada questão.
Em seguida será feita uma análise do desempenho de Língua Portuguesa.
Os comentários da Banca de docentes corretores serão apresentados para o conjunto de questões.
Cumpre esclarecer que, tendo em vista que as questões discursivas de Formação Geral são padronizadas, ou seja, constam de todas as provas, os comentários da Banca são os mesmos para todas as carreiras acadêmicas, sendo direcionados a todos os estudantes que participaram do ENADE/2013.
A seguir, serão analisados os desempenhos da Área de Medicina Veterinária nas duas questões discursivas de Formação Geral do ENADE/2013, comparando os resultados obtidos com comentários para cada questão. 3.3.1.1 Análise de Conteúdo Questão Discursiva 1 do Componente de Formação Geral Os dados de Medicina Veterinária, obtidos a partir das respostas à questão 1, encontram -se na Tabela 3.9 e no Gráfico 3.16.
Nessa questão – de melhor desempenho dentre as duas de Formação Geral – os alunos de todo Brasil tiveram média, 44,1.
A maior média para a questão 1 foi obtida na região Sul ( 45,7 ), e a menor, na região Norte ( 36,3 ).
Quanto à variabilidade das notas, o desvio padrão de todo o Brasil foi 25,0.
O menor desvio padrão foi obtido na região Sul ( 23,7 ) e o maior desvio padrão foi obtido na região Norte ( 27,3 ).
A mediana da questão discursiva 1 foi 50,0 para as regiões Nordeste, Sudeste e Sul, e mais baixa para as regiões Norte ( 40,0 ) e Centro-Oeste ( 45,0 ).
Para o Brasil como um todo foi também 50,0.
As notas máximas ( 100,0 ) foram as mesmas para todas as regiões do Brasil, sem exceção.
Além disso, a nota mínima foi zero em todas as regiões do país, sem exceção.
O Gráfico 3.16 mostra a distribuição das notas na questão discursiva 1 do Componente de Formação Geral.
Observa -se que a maior frequência corresponde aos alunos que obtiveram nota no intervalo ( 40 ; 50 ], correspondendo à moda da distribuição.
A distribuição das notas possui assimetria negativa ( -0,49 ), assim como em quase todas as regiões do Brasil ( entre -0,61 para o Sul e -0,27 para o Centro-Oeste ), excetuando -se apenas a região Norte ( 0,16 ) que possui assimetria positiva. 3.3.1.2 Comentários sobre a correção de Conteúdo das respostas à Questão Discursiva 1 O enunciado restou claro, elaborado com linguagem técnica, estando inserido nas Diretrizes gerais, com nível médio de profundidade, já que abordava tema vinculado à Área de formação da maior parte dos examinandos.
A divisão das perguntas também ficou clara indicando, passo a passo, o caminho de resposta a seguir, com base no texto.
Verificou -se dificuldade quanto à análise dos programas de governo existentes que contribuem para melhoria do saneamento básico.
Nesse ponto, os alunos referiram, com muita frequência, a ações relacionadas à Saúde, tais como : médico da família, ações educativas, ações ambientais, melhoria nas condições de habitação ( sem necessariamente vincular a programas de governo federal ou estaduais ou municipais ).
Os programas citados foram : Programa Saúde da Família ( PSF ), Programa de Aceleração do Crescimento ( PAC ) e Minha Casa Minha Vida e, ainda, o SUS.
Apesar de uma minoria de examinandos ter indicado o PAC do Governo Federal, a maior parte deles apresentou alguma solução para a questão do saneamento básico.
Os estudantes indicaram, com facilidade, a questão da necessidade do tratamento da água como item essencial do saneamento básico.
Muitos concluintes, através da sua experiência local, citaram situações em que as crianças caminhavam de pés no chão perto de rotas de esgoto a céu aberto.
Também criticaram a inércia governamental em não combater as invasões ilícitas, em locais inadequados, sem estrutura de água potável e também de saneamento básico.
Diante das observações preliminares pode ser extraída uma classificação das respostas obtidas que ficou estabelecida em fracas, médias e boas.
As notas consideradas fracas ( notas de 0 a 3,0 ) foram conferidas àqueles que não conseguiram relacionar o saneamento básico com a importância do tratamento de água, os projetos para resolver o problema e nem indicaram soluções para a questão.
As respostas que tiveram média avaliação ( notas de 3,5 a 7,0 ) tocaram, de algum modo, nos pontos indicados, mesmo que de forma insuficiente.
Já as respostas consideradas boas ( notas de 7,5 a 10,0 ) preencheram os requisitos do padrão de resposta, não logrando nota máxima, na maior parte dos casos, por ausência de indicação de programa governamental, mas sugerindo soluções técnicas para o problema.
As respostas mais comuns foram as consideradas medianas e que indicaram, com frequência, que o problema do saneamento básico estaria nas classes de renda mais \textbf{desfavorecida} e nas regiões consideradas como periferias das grandes cidades.
Também foram frequentes respostas que citaram a necessidade de tratamento de água, acompanhada da necessidade de maiores investimentos nessa área. 3.3.1.3 Análise de Conteúdo da Questão Discursiva 2 do Componente de Formação Geral A Tabela 3.10 mostra que o desempenho dos estudantes na questão 2 ( média 39,2 ) foi inferior ao obtido na questão de número 1 ( média 44,1 ).
A região Sul foi aquela onde a média, nessa questão, foi maior ( 40,8 ), e a de menor média foi a região Norte ( 27,9 ).
Quanto à variabilidade das notas, o desvio padrão de todo o Brasil foi 25,8, superior ao obtido na questão de número 1 ( 25,0 ).
O maior desvio nessa questão foi obtido na região Norte ( 28,1 ), enquanto o menor foi obtido na região Sul ( 24,6 ).
A mediana da questão discursiva 2 foi 45,0 para as regiões Sudeste e Sul, e mais baixa para as regiões Norte ( 20,0 ), Nordeste ( 40,0 ) e Centro-Oeste ( 40,0 ).
Para o Brasil como um todo foi também 45,0.
As notas máximas ( 100,0 ) foram as mesmas em todas as regiões do Brasil, assim como as notas mínimas ( 0,0 ), sem exceção.
O Gráfico 3.17 mostra a distribuição das notas na questão discursiva 2 do Componente de Formação Geral.
Neste gráfico, observa -se a grande quantidade de alunos que deixaram a questão 2 em branco ( em torno de 17 \% ), além da moda no intervalo ( 40 ; 50 ].
O coeficiente de assimetria das notas nesta questão é negativo ( -0,21 ), indicando que a distribuição é assimétrica à esquerda.
O mesmo ocorre para quase todas as regiões da Brasil ( entre -0,30 na região Sudeste e -0,02 na Nordeste ), excetuando -se apenas a região Norte ( 0,46 ) que possui assimetria positiva. 3.3.1.4 Comentários sobre a correção de Conteúdo das respostas à Questão Discursiva 2 O tema abordado no enunciado, apesar da complexidade, é extremamente atual.
Essa atualidade facilitou a elaboração de alguma resposta, permitindo que todos os que se dispuseram a responder pudessem apresentar, pelo menos, alguns dos elementos integrantes do padrão de resposta.
A correção, seguindo o padrão de resposta, avaliou a capacidade de elaborar um texto dissertativo que abordasse o tema sob três pontos de vista : o papel da tecnologia digital na atualidade ; a garantia dos direitos do cidadão e do Estado diante do avanço tecnológico ; e o problema da segurança digital para o indivíduo, para o Estado e para as organizações ( empresas e/ou instituições nacionais ou internacionais ).
Dada a amplitude de respostas possíveis, os alunos tiveram desempenho satisfatório.
Identificamos que os concluintes, basicamente, trilharam dois caminhos : a ) a análise da influência das inovações tecnológicas, notadamente através das redes sociais, numa perspectiva essencialmente privada, daí defluindo as questões da privacidade e da segurança ; b ) uma perspectiva de segurança nacional-estatal com criticas à atuação americana de invasão de privacidade dos dados governamentais, e à inercia do governo brasileiro em reagir ao ataque, considerado, pela maioria esmagadora dos concluintes, como indevido e antiético.
Poucos foram os que uniram as duas facetas da questão : o público e o privado.
Na maior parte das provas, o concluinte indicou a ausência de segurança e de privacidade decorrentes das próprias atividades individuais.
Muitos respondentes citaram que, de forma voluntária, os próprios cidadãos enfraquecem a proteção à privacidade, pois divulgam sua vida pessoal, aspectos de seu cotidiano, por meio de imagens e do fornecimento de dados pessoais.
No aspecto segurança, foram inúmeras as indicações de acessos indevidos a dados pessoais com a perpetração de fraudes de todos os tipos, notadamente as bancárias e comerciais, com a aquisição de bens à revelia dos titulares dos dados.
Também no tema segurança, muitos indicaram o aspecto positivo decorrente da vigilância generalizada, não somente através da rede de computadores, mas também pela quantidade de câmeras existentes, que contribuem para a identificação de criminosos procurados ou dos que realizam crimes pela internet.
Na perspectiva pública, a esmagadora maioria indicou repulsa pelo comportamento americano de invasão dos dados do governo brasileiro, reclamando reação do mesmo.
Alguns poucos justificaram a intervenção americana pela “ natural ” disputa de poder pelos Estados imperialistas. 3.3.1.5 Análise de Língua Portuguesa das Questões Discursivas do Componente de Formação Geral Os dados de Medicina Veterinária, obtidos a partir das respostas às questões discursivas do Componente de Formação Geral no que tange à Língua Portuguesa, encontram -se na Tabela 3.11 e no Gráfico 3.18.
Nesse aspecto, os alunos de todo Brasil tiveram média, 48,2.
A maior média com respeito à Língua Portuguesa foi obtida na região Sul ( 50,3 ), e a menor, na região Centro-Oeste ( 45,0 ).
Quanto à variabilidade das notas, o desvio padrão de todo o Brasil foi 21,7.
O menor desvio padrão foi obtido na região Sul ( 19,8 ) e o maior desvio padrão foi obtido na região Norte ( 25,7 ).
A mediana da nota de Língua Portuguesa foi 55,0 para quase todas as regiões do Brasil, exceto as regiões Sudeste e Centro-Oeste onde a mediana foi 52,5.
Para o Brasil como um todo foi também 55,0.
As notas máximas foram as mesmas para quase todas as regiões do Brasil ( 90,0 ), exceto as regiões Sul e Centro-Oeste onde a nota máxima foi 87,5.
Além disso, a nota mínima foi zero em todas as regiões do país, sem exceção.
O Gráfico 3.18 mostra a distribuição das notas de Língua Portuguesa do Componente de Formação Geral.
Observa -se que a maior frequência corresponde aos alunos que obtiveram nota no intervalo ( 50 ; 60 ], correspondendo à moda da distribuição.
A distribuição das notas possui assimetria negativa ( -1,05 ), assim como em todas as regiões do Brasil ( entre -1,23 para o Sul e -0,78 para o Norte ). 3.3.1.6 Comentários sobre a correção das respostas de Formação Geral com respeito à Língua Portuguesa Um aspecto interessante a ser destacado é a diferença de desempenho entre a questão 1 e a questão 2, principalmente nos aspectos textuais – os textos produzidos para a questão 1 apresentaram resultados melhores do que os produzidos para a questão 2.
Uma hipótese que provavelmente explique essa discrepância é o fato de o tema “ saneamento básico ” ser mais relevante para profissionais de saúde do que o tema “ segurança e privacidade ”, constituindo tema trabalhado no currículo escolar.
Prova disso é que quase todos os textos da questão 1 citam termos específicos da área de saúde, como “ contaminação feco-oral ”, “ higiene corporal ”, “ ingestão de alimentos ou água contaminada por vetores biológicos ”.
Os resultados da avaliação correspondem aos seguintes aspectos observados em cada competência : a ) aspectos ortográficos : o desempenho dos participantes ficou concentrado, majoritariamente, nos níveis 3 ( até 3 desvios ) e 2 ( de 4 a 7 desvios ), devido ao grande índice de desvios de acentuação, contrabalançada pelo pequeno número de desvios em grafia.
Foi reduzido o número de textos com ausência de desvios ( enquadrados no nível 4 ).
Alguns textos, com pior desempenho, foram enquadrados no nível 1 ( de 8 a 12 desvios ), enquanto o nível zero foi atribuído a pouquíssimos casos, já que ele revela falta absoluta de domínio das convenções ortográficas.
Observou -se, portanto, que existe uma diferença muito grande de desempenho nos dois aspectos analisados : baixíssimo índice de desvios ortográficos e grande índice de desvios de acentuação.
Em vários casos, ocorre ausência completa de acentuação gráfica.
Os resultados revelam, portanto, que a tendência dominante entre os universitários brasileiros é a eliminação da acentuação gráfica, talvez motivada pela ausência de esclarecimento dos meios de comunicação, das autoridades e das escolas sobre as decisões do Acordo Ortográfico da Língua Portuguesa de 1990.
Os casos mais sistemáticos de ausência de acentuação são :  palavras proparoxítonas ( “ politicas ”, “ publicas ”, “ juridicas ” ) ;  palavras paroxítonas terminadas em ditongo crescente ( “ varios ”, “ Brasilia ”, “ familia ” ) ;  palavras oxítonas ( “ ninguem ”, “ esta ”, “ ate ” ).
Quanto ao domínio das convenções relativas à grafia das palavras, são pontuais os problemas encontrados ( “ infecciozas ”, “ obtos ”, “ hipedemia ”, “ saniamento ”, “ cituação ”, “ abitantes ” ).
Os únicos desvios recorrentes foram : a ausência de segmentação nos casos de verbo + pronome oblíquo “ se ”, combinada, eventualmente, à alteração da grafia do pronome ( “ percebese ”, “ percebece ”, “ preferece ” ) ; a grafia da forma verbal de pretérito imperfeito do subjuntivo como se fosse uma ocorrência de verbo + pronome oblíquo ( “ \textbf{distribui} -se ” por “ \textbf{distribuísse} ” ) Vale observar, também, que, ao contrário do que se esperava, não apareceram abreviaturas próprias do “ internetês ”, ou seja, dos hábitos de comunicação escrita adquiridos pelo uso de redes sociais e e-mails. b ) aspectos textuais : o desempenho dos participantes ficou concentrado, majoritariamente, nos níveis 3 e 2, devido à grande ocorrência de problemas de estruturação textual.
Foi muito baixo o número de textos enquadrados no nível 4 por não apresentar qualquer problema estrutural.
Alguns textos com pior desempenho foram enquadrados no nível 1, enquanto o nível zero foi atribuído a pouquíssimos casos, já que ele revela a existência de um texto sem articulação e com comprometimento do sentido.
Observou -se que a grande maioria dos participantes não \textbf{distribuiu} as ideias em parágrafos, talvez devido ao pequeno número de linhas disponibilizadas para a resposta da questão ou, quem sabe, pela suposição de que não seria necessária essa divisão por não se tratar de um texto no modelo de uma redação dissertativo-argumentativa, como solicitado nos vestibulares.
Essa competência é a que se revela como a mais problemática entre os participantes, porque são muitos os problemas observados, desvios acumulados durante toda a formação do estudante e que não se resolvem com um estudo autodidata, como acontece com regras ortográficas ou morfossintáticas : sequência justaposta de ideias sem encaixamentos sintáticos ; redução drástica de estruturas subordinadas, ao lado do aumento na frequência de estruturas coordenadas e absolutas ; redução no uso de conectores para expressar relações lógicas essenciais à construção do texto, substituídas pela exigência de inferência por parte do interlocutor para suprir a sua ausência ; emprego equivocado de operadores que não estabelecem relações lógicas coerentes entre ideias do texto ; emprego inadequado do pronome relativo ( com omissão da preposição ou a utilização de pronome inadequado, como “ onde ” ) ; repetição exaustiva de termos sem a utilização de procedimentos mais sofisticados de substituição ( hiperonímias, hiponímias, nominalizações, expressões metafóricas ) ; frases fragmentadas que comprometem a estrutura lógico-gramatical ; frases formadas apenas por oração subordinada, sem oração principal.
Um importante aspecto a destacar é o baixíssimo desempenho de uma parte dos participantes em relação à estrutura formal do texto produzido, o que é extremamente preocupante ao se levar em conta que são graduandos em fase final de formação.
São frequentes os casos de desvios de estruturação frasal, com uso inadequado ou ausência de conectivos entre parágrafos e entre frases.
Em uma parte dos textos, falta um mínimo de textualidade e de domínio do registro padrão da língua.
Na verdade, observam -se relações linguísticas quase agramaticais, como as estabelecidas pela sequência de gerúndios sem o apoio de um ponto de partida para a organização das informações gramaticais e semânticas.
Quanto à utilização dos mecanismos de referenciação, dois aspectos devem ser destacados : a ocorrência, em uma boa parte dos textos, de repetições de palavras ou expressões sem a utilização de termos sinônimos ou pronomes, como seria adequado ; e a alternância entre os pronomes demonstrativos este / esse, e suas formas flexionadas, sem a observância de um critério de utilização.
Esse último aspecto tem se generalizado na modalidade oral da língua e, cada vez mais, tem se tornado uma constante na modalidade escrita de jornais e revistas : predominância da forma “ esse ” na modalidade oral e da forma “ este ” na modalidade escrita, provavelmente por um processo de hipercorreção. c ) aspectos morfossintáticos e vocabulares : o desempenho dos participantes ficou concentrado, majoritariamente, nos níveis 3 e 2, devido à grande ocorrência de problemas de regência, concordância e uso adequado dos sinais de pontuação.
O nível 4 foi atribuído a um número muito reduzido de textos.
Alguns textos com pior desempenho foram enquadrados no nível 1, enquanto o nível zero foi atribuído a pouquíssimos casos, já que ele revela a existência de um texto sem articulação e com comprometimento do sentido.
Os resultados são muito transparentes em relação aos aspectos mais problemáticos no desempenho dos participantes.
O desvio mais frequente é a falta do sinal indicativo da crase, o que se configura como um problema de regência, verbal ou nominal, conforme o caso – isso revela que o usuário não tem consciência de que, sob a forma do termo “ a ” existe a presença de uma preposição “ a ”, exigida pela regência do termo anterior.
Embora em outros exames, como o Enem, a falta de crase seja apenada em acentuação, nesta correção amostral esse desvio foi considerado no âmbito dos aspectos morfossintáticos.
Outro problema relacionado à regência verbal e à nominal, encontrado frequentemente nas questões, foi a ausência de preposição antes de pronome relativo, processo generalizado na modalidade oral da língua, em situações de registro informal.
Apesar da possibilidade de que essa alteração de regência se generalize no padrão escrito da Língua Portuguesa, como já está ocorrendo até em textos jornalísticos, o não emprego da preposição foi apenado neste processo de avaliação. preocupante ao se levar em conta que são graduandos em fase final de formação.
São frequentes os casos de desvios de estruturação frasal, com uso inadequado ou ausência de conectivos entre parágrafos e entre frases.
Em uma parte dos textos, falta um mínimo de textualidade e de domínio do registro padrão da língua.
Na verdade, observam -se relações linguísticas quase agramaticais, como as estabelecidas pela sequência de gerúndios sem o apoio de um ponto de partida para a organização das informações gramaticais e semânticas.
Quanto à utilização dos mecanismos de referenciação, dois aspectos devem ser destacados : a ocorrência, em uma boa parte dos textos, de repetições de palavras ou expressões sem a utilização de termos sinônimos ou pronomes, como seria adequado ; e a alternância entre os pronomes demonstrativos este / esse, e suas formas flexionadas, sem a observância de um critério de utilização.
Esse último aspecto tem se generalizado na modalidade oral da língua e, cada vez mais, tem se tornado uma constante na modalidade escrita de jornais e revistas : predominância da forma “ esse ” na modalidade oral e da forma “ este ” na modalidade escrita, provavelmente por um processo de hipercorreção. c ) aspectos morfossintáticos e vocabulares : o desempenho dos participantes ficou concentrado, majoritariamente, nos níveis 3 e 2, devido à grande ocorrência de problemas de regência, concordância e uso adequado dos sinais de pontuação.
O nível 4 foi atribuído a um número muito reduzido de textos.
Alguns textos com pior desempenho foram enquadrados no nível 1, enquanto o nível zero foi atribuído a pouquíssimos casos, já que ele revela a existência de um texto sem articulação e com comprometimento do sentido.
Os resultados são muito transparentes em relação aos aspectos mais problemáticos no desempenho dos participantes.
O desvio mais frequente é a falta do sinal indicativo da crase, o que se configura como um problema de regência, verbal ou nominal, conforme o caso – isso revela que o usuário não tem consciência de que, sob a forma do termo “ a ” existe a presença de uma preposição “ a ”, exigida pela regência do termo anterior.
Embora em outros exames, como o Enem, a falta de crase seja apenada em acentuação, nesta correção amostral esse desvio foi considerado no âmbito dos aspectos morfossintáticos.
Outro problema relacionado à regência verbal e à nominal, encontrado frequentemente nas questões, foi a ausência de preposição antes de pronome relativo, processo generalizado na modalidade oral da língua, em situações de registro informal.
Apesar da possibilidade de que essa alteração de regência se generalize no padrão escrito da Língua Portuguesa, como já está ocorrendo até em textos jornalísticos, o não emprego da preposição foi apenado neste processo de avaliação.
Outro desvio muito frequente diz respeito aos processos de concordância verbal e de concordância nominal.
Quanto à concordância de número, observou -se ausência de marca ( com sujeito anteposto ou posposto ) ou uso indevido ( uso inadequado da marca de plural comandado pelo núcleo plural da locução adjetiva, apesar de o substantivo que funciona como núcleo do sintagma nominal estar no singular ).
Uma ocorrência que se destacou foi a ausência de acento circunflexo na forma plural do presente do indicativo do verbo “ ter ”, que foi considerada como um desvio na concordância verbal e não na acentuação gráfica.
Quanto à concordância de gênero, vários casos foram observados, normalmente no âmbito de sintagmas nominais longos, em que o adjetivo está afastado do substantivo.
Quanto à questão da colocação pronominal, foram poucos os casos observados.
Apesar de serem aspectos relacionados à oralidade, concluiu -se que, no registro escrito formal, a maioria dos participantes já incorporou regras como a não introdução da frase por um pronome oblíquo e a próclise na presença de um termo atrator.
Não se adotou, entretanto, o padrão excessivamente formal descrito pelas gramáticas normativas em relação à posição do pronome oblíquo em locuções verbais, já que esse uso está muito distante da prática cotidiana, até em textos mais formais.
Quanto aos aspectos vocabulares, alguns tipos de inadequação foram observados : expressões da oralidade apareceram em algumas respostas, mas sem maior relevância do ponto de vista quantitativo ; seleção vocabular incompatível com o contexto, gerando situações de falta de inteligibilidade ; falta de domínio de vocabulário mais abstrato e de maior complexidade, essencial ao desenvolvimento do texto de base dissertativa.
Quanto à utilização dos sinais de pontuação, observou -se uma grande precariedade nos textos analisados.
Não foi apenada a ausência de vírgula para destacar locuções ou adjuntos adverbiais de pequena extensão deslocados de posição na frase, por ser um uso opcional.
São os seguintes os tipos de problemas encontrados : a ) vírgula : utilização de vírgula para separar o sujeito e o predicado ; ocorrência de apenas uma das vírgulas para separar uma palavra, uma expressão ou uma oração encaixada ; uso de vírgula no lugar do ponto para separar ideias que constituem períodos distintos ; ausência de vírgula para separar enumerações ; ausência de vírgula para separar oração adjetiva explicativa ou utilização inadequada para separar oração adjetiva restritiva ; b ) ponto e vírgula : utilização do ponto e vírgula no lugar de vírgula ; c ) ponto final : ausência de ponto final para separar períodos.
Várias marcas de oralidade foram identificadas, embora não em alta frequência : o uso do pronome relativo “ onde ” como relativo universal, falta de artigo definido antes de 3.3.2 Componente de Conhecimento Específico Na parte da prova relativa às questões discursivas no Componente de Conhecimento Específico ( Tabela 3.12 ), observa -se que a média foi mais baixa do que nas questões discursivas do Componente de Formação Geral.
Enquanto no Componente de Formação Geral a média para estudantes de Medicina Veterinária de todo o Brasil foi 42,9, na Parte de Conhecimento Específico a média foi 18,8.
A maior média deste componente foi obtida pelos estudantes da região Sul ( 20,6 ), e a menor, pelos da região Norte ( 12,4 ).
Quanto à variabilidade das notas, o desvio padrão de todo o Brasil foi 14,0.
O maior desvio padrão foi encontrado na região Sul ( 14,4 ), e o menor, na região Norte ( 10,6 ).
A maior nota máxima, nas questões discursivas do Componente de Conhecimento Específico, foi obtida na região Sudeste ( 73,3 ), enquanto a menor nota máxima foi encontrada na região Norte ( 45,0 ).
Além disso, a nota mínima ( 0,0 ) foi obtida por alunos de todas as regiões do Brasil sem exceção.
A mediana do Brasil como um todo foi 16,7, também encontrada na região Centro-Oeste.
As demais medianas foram : 10,0 na região Norte, 15,0 na região Nordeste, 18,3 na região Sudeste e 20,0 na região Sul.
O Gráfico 3.19 representa a distribuição das notas nas questões discursivas no Componente de Conhecimento Específico.
Essa distribuição tem moda no primeiro intervalo, o [ 0 ; 10 ].
Nota -se que há uma tendência decrescente a partir deste intervalo até o antepenúltimo intervalo.
Nenhum aluno recebeu notas alocadas nos dois últimos intervalos.
O coeficiente de assimetria é positivo para todo o Brasil ( 0,62 ) e para todas as regiões ( entre 0,56 para as regiões Sudeste e Sul, e 0,78 para a região Centro-Oeste ), indicando uma assimetria à direita. 3.3.2.1 Análise da Questão Discursiva 3 do Componente de Conhecimento Específico Na questão 3, cujos resultados aferidos encontram -se descritos na Tabela 3.13, a média dos estudantes de todo o Brasil foi 15,8.
A menor média nessa questão foi obtida pelos alunos da região Norte ( 9,7 ), enquanto a maior média foi obtida na região Sul ( 18,9 ).
Quanto à variabilidade das notas, o desvio padrão de todo o Brasil foi 21,3.
O maior desvio padrão foi obtido na região Sudeste ( 22,5 ), enquanto o menor foi obtido nas regiões Norte e Centro-Oeste ( 16,5 ).
A maior nota máxima foi alcançada por pelo menos um aluno na região Sul ( 100,0 ), enquanto a menor nota máxima foi obtida na região Norte ( 75,0 ).
A mediana e a nota mínima do Brasil como um todo foi zero, o mesmo valor para quase todas as regiões, exceto a mediana da região Sul ( 10,0 ), indicando que pelo menos metade dos estudantes de cada região recebeu nota zero nesta questão, seja por terem deixado a questão em branco, seja pela inadequação da resposta.
O Gráfico 3.20 mostra a distribuição das notas na questão discursiva 3 do Componente de Conhecimento Específico.
Essa distribuição apresenta uma moda nas questões com nota zero.
O coeficiente de assimetria é positivo para todo o Brasil ( 1,15 ) e para todas as regiões ( entre 0,85 na região Sul e 1,91 na região Norte ), indicando uma assimetria à direita. 3.3.2.2 Comentários sobre as respostas à Questão Discursiva 3 O enunciado da questão era claro, indicando que nas respostas os alunos considerassem apenas os impactos socioeconômicos e de segurança alimentar, ao solicitar :"Considerando as atribuições do Sistema Brasileiro de Inspeção ( SISBI ), redija um texto dissertativo sobre as implicações socioeconômicas (... ) oriundos da agricultura familiar no SISBI".
Porém, o padrão de resposta exigia mais do que isso.
O contexto utilizado na questão era bastante pertinente, atual e previsto nos conteúdos programáticos de disciplinas do curso de Medicina Veterinária.
As respostas dadas pelos estudantes, em sua maioria, ficaram restritas aos impactos socioeconômicos e de segurança alimentar, como induzia o enunciado da questão.
Houve respostas indicando soluções divergentes do esperado pelo padrão : alguns alunos listaram a dificuldade financeira pela qual passariam os agricultores familiares em decorrência da implantação do SISBI ( impacto socioeconômico ) e escreveram sobre o risco de comercialização de produtos de agricultura familiar, o que poderia aumentar o número de pessoas infectadas por patógenos veiculados por alimentos oriundos da produção animal.
Este fato chamou atenção, pois o que os alunos demonstraram foi desconhecimento sobre o SISBI.
A questão envolvia a redação de um texto envolvendo uma análise subjetiva e os alunos demonstrando dificuldade de compreensão do que foi solicitado.
Eles responderam de forma truncada e evidenciaram dificuldade de redação e de síntese ( que talvez seja um reflexo do panorama nacional da educação ).
Na maioria das respostas, os alunos atribuíram um título, acreditando que, uma vez que a questão solicitou um texto, deveriam seguir uma estrutura de redação escolar, mesmo que o texto fosse constituído de um só período ou parágrafo.
Quanto ao conteúdo, a grande maioria das respostas foi fraca, em especial porque quatro pontos deveriam ser atribuídos à dissertação sobre o panorama atual do Sistema de Inspeção e o papel do MAPA frente à implantação do SISB, o que estava pouco presente nas respostas dos alunos.
A maior parte das respostas dadas baseou -se apenas no impacto socioeconômico e de segurança alimentar.
A maioria dos alunos redigiu sobre o aumento na renda e ampliação de mercado para o produtor familiar, assim como sobre a referência de um produto seguro para a população, em virtude de um programa para inspeção de produtos oriundos da agricultura familiar.
Entre as respostas apresentadas pelos alunos, a referência sobre os impactos sociais na implantação do SISBI foi a mais frequente.
Além disso, verificou -se que a maioria dos alunos mencionava exclusivamente a participação do médico veterinário no Sistema de Inspeção e a necessidade do MAPA abrir concurso e vagas para a prática da inspeção.
Foram poucas as respostas boas.
Notou -se que a interpretação proposta pelo padrão de resposta para a questão só foi atendida por um grupo restrito de alunos, que conseguiu abordar a situação atual da inspeção, bem como o papel do MAPA na implantação do SISBI.
A média baixa justifica -se pela apresentação de respostas que abordavam somente os impactos socioeconômicos ou somente de segurança alimentar, como já comentado.
O número de respostas com valores zerados foi considerável, o que pode ser justificado pela dificuldade de interpretação da questão e/ou por desconhecimento sobre o SISBI.
Em grande parte das respostas, os alunos indicavam a necessidade de apoio e serviços, que já haviam sido incluídos para implantação do SISBI, e as respostas estavam totalmente fora do contexto da questão.
Algumas notas zero foram atribuídas a respostas consideradas anuladas por conterem conteúdo agressivo, palavras de baixo calão, protesto político.
Existe, por parte dos alunos, grande dificuldade na expressão ortográfica e gramatical.
Boa parte dos textos redigidos continha erros ortográficos e de concordância verbal e nominal, demonstrando necessidade de atenção geral ao Sistema de Educação Nacional. 3.3.2.3 Análise da Questão Discursiva 4 do Componente de Conhecimento Específico A Tabela 3.14 contém as informações relativas à questão 4 do conjunto de questões do Componente de Conhecimento Específico.
O desempenho médio dos estudantes de todo o Brasil nesta questão foi inferior ao desempenho nas questões de números 3 e 5.
A média geral do Brasil foi 14,3, sendo a menor média registrada na região Norte ( 5,8 ) e a maior na região Sudeste ( 16,7 ).
Quanto à variabilidade das notas, o desvio padrão de todo o Brasil foi 17,9.
O maior desvio padrão foi obtido na região Sudeste ( 18,4 ), enquanto o menor foi obtido na região Norte ( 12,7 ).
A nota máxima ( 100,0 ) foi atingida em quase todas as regiões, exceto na região Sudeste ( 90,0 ).
A mediana e a nota mínima foram zero para as regiões Norte, Nordeste e Centro-Oeste, indicando que, nestas regiões, pelo menos metade dos alunos tirou nota zero nesta questão.
A mediana do Brasil como um todo foi 5,0, sendo também encontrada na região Sul, enquanto na região Sudeste a mediana foi 15,0.
A nota mínima foi zero para todas as regiões do Brasil, sem exceção.
O Gráfico 3.21 representa a distribuição de notas na questão discursiva 4, no Componente de Conhecimento Específico.
O número de estudantes que deixou a questão em branco foi muito grande, correspondendo à moda da distribuição.
O coeficiente de assimetria é positivo para todo o Brasil ( 1,30 ) e para todas as regiões ( entre 1,02 para a região Sudeste e 3,02 para a região Norte ), indicando uma assimetria acentuada à direita. 3.3.2.4 Comentários sobre as respostas à Questão Discursiva 4 A Questão 4 compreendia três itens, que exploravam a capacidade de aplicação de conhecimentos adquiridos em disciplinas como Anatomia, Fisiologia, Farmacologia, Diagnóstico por Imagem, Patologia Geral e Clínica Médica, na forma de discussão de um caso clínico.
Não houve muita divergência entre o que se esperava no padrão de resposta e as respostas dos alunos para a avaliação dos itens ‘ a ’ e ‘ c ’.
Em relação ao item ‘ b ’, a pergunta não deixou clara de que maneira a enfermidade deveria ser “ identificada ” e, por isso, o padrão de resposta poderia ser mais flexível.
Um grande número de estudantes descreveu a enfermidade como sendo Insuficiência Cardíaca Congestiva ou Cardiomiopatia Dilatada Esquerda, e explicavam a patogenia das situações especificadas, o que não foi previsto pelo padrão de resposta.
O item ‘ a ’ foi respondido de forma bem objetiva, pela maioria dos participantes, sendo reduzido o número de respostas em branco.
Os estudantes apresentaram um desempenho mediano para responder ao item ‘ a ’.
Muitas vezes, apesar de corretas as interpretações da imagem ecocardiográfica apresentada na questão, faltava o uso adequado de termos técnicos, sendo inclusive observado em uma prova o diagnóstico"hipermetropia de parede cardíaca", termo utilizado em oftalmologia.
Esse item deixou claro que há bastante dificuldade de interpretação da imagem ecocardiográfica.
Houve respostas que chamaram atenção : muitos relataram a ausência de ventrículo direito ; alguns se referiram a fígado, pulmão, à esplenomegalia do ventrículo esquerdo.
O item ‘ b ’, que solicitava que o estudante indicasse “ a patogenia da enfermidade diagnosticada ” a partir do caso clínico, apesar da clareza do enunciado, apresentou o menor índice de acerto e o maior índice de respostas em branco.
O diagnóstico de Insuficiência Cardíaca Congestiva e/ou cardiomiopatia Hipertrófica/Dilatada foi muito frequente.
Provavelmente, esse fato ocorreu por erro interpretativo do que foi questionado, demonstrando que muitos alunos confundem patogenia com patologia ( relacionada ao diagnóstico da enfermidade no paciente em questão ).
As respostas dos alunos evidenciaram, também, problemas em relação aos seguintes aspectos : - Falta de diagnóstico, pois não utilizaram as alterações clínicas apresentadas no enunciado e as alterações apresentadas na imagem para o estabelecimento adequado da enfermidade ; - Erro no estabelecimento do diagnóstico da enfermidade.
A grande dificuldade em responder o que era solicitado no item ‘ b ’ parece ser fruto de falta de conhecimento científico, como também desconhecimento de termos técnicos apropriados.
Algumas respostas podem ser consideradas absurdas : hepatopatia congênita na imagem ecocardiográfica ; pneumonia ; úlcera ; filariose.
Tal dificuldade se justifica, também, pelo item exigir conhecimento mais detalhado das disciplinas relacionadas ao estudo do caso apresentado, sendo o item com o maior grau de complexidade e profundidade.
O item ‘ c ’ solicitava a indicação da terapia medicamentosa, mas não deixava claro se deveriam ser utilizados nomes comerciais, nomes químicos, nomes genéricos ou grupos farmacológicos dos medicamentos.
Dessa forma, foi grande o número de respostas diferentes, algumas previstas pelo padrão de resposta, outras não.
Também foi observada uma espécie de memorização dos estudantes para os tratamentos dos problemas cardíacos, pois em várias provas as respostas apresentavam apenas siglas como, por exemplo : tratamento dos 4D’s - digitálicos, diuréticos, descanso e dieta.
Esse item, apesar de também ter sido, em algumas situações, composto por respostas que demonstram uma grande falta de conhecimento e de capacidade de raciocínio ( como o uso de medicamentos 3.3.2.5 Análise da Questão Discursiva 5 do Componente de Conhecimento Específico A Tabela 3.15 contém as informações relativas à questão 5 do conjunto do Componente de Conhecimento Específico.
O desempenho dos estudantes nessa questão foi superior ao das questões 3 e 4.
A nota média dos estudantes de todo o Brasil foi 26,3.
A maior média foi registrada na região Centro-Oeste ( 28,8 ), enquanto a menor média foi registrada na região Norte ( 21,7 ).
Quanto à variabilidade das notas, o desvio padrão dos alunos do Brasil, como um todo, foi 24,3.
Enquanto o maior desvio foi encontrado na região Centro-Oeste ( 24,7 ), o menor foi encontrado na região Norte ( 21,3 ).
A nota máxima ( 100,0 ) foi alcançada em todas as regiões do Brasil, exceto na região Norte ( 80,0 ).
A mediana foi 20,0 para quase todas as regiões do Brasil, exceto para as regiões Sul ( 30,0 ) e Centro-Oeste ( 25,0 ).
A nota mínima foi zero para todas as regiões do Brasil, sem exceção. hipertensivos para pacientes com insuficiência cardíaca, medicamentos anti-diuréticos para pacientes com predisposição ao desenvolvimento de edema, medicamentos “ cardiopatas ” ), apresentou um índice de acertos mais elevado que os demais itens.
Se por um lado foram poucas as respostas com protestos contra o ENADE e/ou à Instituição de Ensino do estudante, por outro, a quantidade de erros de português presente nas respostas redigidas pelos alunos foi assustadora.
Houve, inclusive, erros na ortografia de termos frequentes no cotidiano de alunos da Área da Saúde.
O número de respostas em branco foi considerável e pode estar relacionado com a dificuldade de interpretação da imagem, visto que nem todas as IES possuem aparelhos para a realização de ecocardiograma, e no geral há pouco estudo prático sobre o assunto. 3.3.2.5 Análise da Questão Discursiva 5 do Componente de Conhecimento Específico A Tabela 3.15 contém as informações relativas à questão 5 do conjunto do Componente de Conhecimento Específico.
O desempenho dos estudantes nessa questão foi superior ao das questões 3 e 4.
A nota média dos estudantes de todo o Brasil foi 26,3.
A maior média foi registrada na região Centro-Oeste ( 28,8 ), enquanto a menor média foi registrada na região Norte ( 21,7 ).
Quanto à variabilidade das notas, o desvio padrão dos alunos do Brasil, como um todo, foi 24,3.
Enquanto o maior desvio foi encontrado na região Centro-Oeste ( 24,7 ), o menor foi encontrado na região Norte ( 21,3 ).
A nota máxima ( 100,0 ) foi alcançada em todas as regiões do Brasil, exceto na região Norte ( 80,0 ).
A mediana foi 20,0 para quase todas as regiões do Brasil, exceto para as regiões Sul ( 30,0 ) e Centro-Oeste ( 25,0 ).
A nota mínima foi zero para todas as regiões do Brasil, sem exceção.
O Gráfico 3.22 representa a distribuição das notas da questão discursiva 5 do Componente de Conhecimento Específico.
Esta distribuição apresenta várias modas, sendo a moda principal no intervalo ( 0 ; 10 ].
Destaca -se, também, o intervalo ( 30 ; 40 ] que se apresenta como uma segunda moda.
Entre as três questões discursivas do Componente de Conhecimento Específico, esta foi a que apresentou um maior espalhamento das notas daqueles que resolveram a questão.
O coeficiente de assimetria para todos os alunos é 0,69 e para todas as regiões se mantém positivo, variando entre 0,55 na região Centro-Oeste e 0,77 na região Nordeste. 3.3.2.6 Comentários sobre as respostas à Questão Discursiva 5 A questão foi bem elaborada, com enunciado claro e de fácil compreensão.
Abordava o tema"confinamento"que além de pertinente tem abrangência nacional.
Além disso, a questão trata de assunto que tem sido divulgado e discutido pela mídia, em especial a televisiva.
O padrão de resposta sugerido pela Comissão Assessora abrangia bem as possibilidades de resposta, a menos do item ‘ a ’, para o qual existem outras vantagens do confinamento que não foram consideradas.
Em relação ao item ‘ a ’, de uma forma geral, as respostas dos alunos citavam alguma ou algumas das quatro vantagens presentes no padrão determinado.
Foi possível observar forte influência de temas debatidos em outras disciplinas, o que pode justificar o aparecimento, num conjunto significativo de provas, da resposta"diminui as chances de intoxicação por plantas tóxicas".
Outro tipo de resposta, que apareceu com frequência ao longo de toda a correção, referia -se ao confinamento como forma de redução do desmatamento.
Apesar de não ser uma das respostas válidas, denota uma preocupação com o ambiente.
Para responder ao item ‘ b ’, “ como proceder para reduzir custos de produção em um sistema de produção de bovinos de corte em confinamento ”, bastava tomar por base o quadro apresentado no enunciado.
As respostas consideradas fracas foram as que não continham as vantagens fixadas no padrão de resposta, ou apresentavam apenas alguma, e não respondiam ao item ‘ b ’ corretamente, com base no quadro apresentado.
As respostas medianas ficaram caracterizadas por conterem todas as quatro vantagens fixadas no padrão de resposta e uma das duas principais medidas para a redução do custo operacional ; ou por serem parciais ou no item ‘ a ’ ou no item ‘ b ’.
As vantagens mais citadas pelos alunos foram : possibilidade de detecção de doenças e de estro com mais facilidade ; manejo facilitado ; e necessidade de mão de obra reduzida.
Em relação às medidas de redução de custos, foram comumente relatados como importantes : o controle sanitário ; a oferta de alimentos adequados ; contratação de mão de obra especializada.
Alguns abordaram o uso de alimentos"alternativos", mas sem identificar as vantagens para a redução do custo operacional.
Em algumas provas observou -se um nível alto de desconhecimento sobre o assunto.
Foi bastante comum o encontro de caligrafias ininteligíveis e erros grosseiros de ortografia e construção de frases.
Em algumas respostas, os alunos demonstravam noções de administração rural e de empreendedorismo, em outras, ficava claro que a formação era puramente de clínica médica, voltada para pequenos animais.
O número de respostas em branco foi pequeno, provavelmente por se tratar de uma questão fácil com assunto de abrangência nacional, discutido dentro e fora do curso de graduação.
Apesar disso, a média foi baixa, provavelmente pelo fato da Comissão exigir que os alunos apresentassem as quatro vantagens por eles listadas, apesar da pergunta sequer estabelecer o quantitativo de vantagens que deveriam constar da resposta.
Além disso, como existem muitas vantagens, também não ficou claro que o aluno deveria escolher aquelas mais importantes, de maior relevância.
É provável que a falta de atenção na leitura da questão seja a principal causa dos erros do item ‘ b ’.
Algumas respostas foram consideradas nulas ( também receberam nota zero ) por conterem expressão de desconhecimento ("não sei responder"), protesto contra o ENADE, ou, ainda, opinião contrária à ingestão de carne e criação de animais.

% matched lemmas: desfavorecer, distribuir
\end{textsample}
